\section*{Bab \ref{ch:g6:wholes} --- Bilangan Cacah}

\begin{solution}{exo:g6:wholes:1}
$3\,047$; \quad $20\,500$; \quad $104$. \quad
$60\,013$: ``enam puluh ribu tiga belas''.
\end{solution}

\begin{solution}{exo:g6:wholes:2}
Angka ratusan: $2$; angka satuan: $9$; angka $4$ membilang puluhan
ribu. Uraiannya:
\[
47\,209 = 4 \times 10\,000 + 7 \times 1000 + 2 \times 100 + 0 \times 10
+ 9 .
\]
\end{solution}

\begin{solution}{exo:g6:wholes:3}
$507 < 570$ (puluhan: $0 < 7$); \quad
$1\,099 > 989$ (empat angka mengalahkan tiga); \quad
$23\,456 < 23\,465$ (puluhan: $5 < 6$); \quad
$9\,999 < 10\,000$ (empat angka melawan lima).
\end{solution}

\begin{solution}{exo:g6:wholes:4}
Menurut banyaknya angka, $978$ dan $1\,001$ berada di depan ketiga
bilangan empat angka yang diawali $2$; di antara ketiganya,
bandingkan puluhan dan satuannya: $2\,034 < 2\,043 < 2\,304$. Urutan akhirnya:
\[
978 < 1\,001 < 2\,034 < 2\,043 < 2\,304 .
\]
\end{solution}

\begin{solution}{exo:g6:wholes:5}
$348 + 576 = 924$; \quad $805 - 268 = 537$ (periksa:
$537 + 268 = 805$); \quad $307 \times 26 = 307 \times 20 + 307 \times 6
= 6\,140 + 1\,842 = 7\,982$.
\end{solution}

\begin{solution}{exo:g6:wholes:6}
$38 + 45 + 12 + 5 = (38 + 12) + (45 + 5) = 50 + 50 = 100$.

$2 \times 83 \times 50 = (2 \times 50) \times 83 = 100 \times 83 =
8\,300$.

$4 \times 78 \times 25 = (4 \times 25) \times 78 = 100 \times 78 =
7\,800$.
\end{solution}

\begin{solution}{exo:g6:wholes:7}
$95 = 7 \times 13 + 4$ (\omterm{def:g3:division:remainder}{hasil bagi} $13$, \omterm{def:g3:division:remainder}{sisa} $4$).

$230 = 6 \times 38 + 2$ (\omterm{def:g3:division:remainder}{hasil bagi} $38$, \omterm{def:g3:division:remainder}{sisa} $2$).

$504 = 8 \times 63 + 0$ (\omterm{def:g3:division:remainder}{hasil bagi} $63$, \omterm{def:g3:division:remainder}{sisa} $0$: $504$ \omterm{def:g6:wholes:divisible}{habis
dibagi} $8$).
\end{solution}

\begin{solution}{exo:g6:wholes:8}
Kesamaannya benar, tetapi \omterm{def:g3:division:remainder}{sisa} $13$ \emph{tidak} lebih kecil daripada
\omterm{def:g5:division:multiple}{pembagi} $8$, jadi itu bukan pembagian Euklides. Keluarkan satu $8$
lagi dari sisanya: $173 = 8 \times 21 + 5$, dengan $0 \leq 5 < 8$:
\omterm{def:g3:division:remainder}{hasil bagi} $21$, \omterm{def:g3:division:remainder}{sisa} $5$.
\end{solution}

\begin{solution}{exo:g6:wholes:9}
\emph{1.} $2\,000 = 12 \times 166 + 8$: peternakan itu mengisi $166$
kotak penuh dan $8$ telur tersisa.

\emph{2.} $150 = 12 \times 12 + 6$: dua belas kotak hanya memuat $144$
telur, belum cukup. Toko roti itu harus membeli $13$ kotak.
\end{solution}

\begin{solution}{exo:g6:wholes:10}
$100 = 7 \times 14 + 2$: seratus hari adalah $14$ minggu penuh
ditambah $2$ hari. Empat belas minggu kemudian hari Selasa lagi; dua
hari sesudah Selasa adalah hari \emph{Kamis}.
\end{solution}

\begin{solution}{exo:g6:wholes:11}
\omterm{def:g3:division:remainder}{Sisa} terbesar yang mungkin pada pembagian oleh $9$ adalah $8$ (\omterm{def:g3:division:remainder}{sisanya}
harus lebih kecil daripada \omterm{def:g5:division:multiple}{pembaginya}). Jadi bilangan yang dibagi itu
\[
9 \times 56 + 8 = 504 + 8 = 512 .
\]
\end{solution}

\begin{solution}{exo:g6:wholes:12}
Misalkan angkanya $h$ (ratusan), $0$ (puluhan), $u$ (satuan). Kita
tahu $h = 2 \times u$ dan $h + 0 + u = 9$, jadi $2u + u = 9$, yang
memberi $3u = 9$, $u = 3$, dan $h = 6$. Bilangannya $603$. Periksa:
$6$ adalah dua kali $3$, angka puluhannya $0$, dan $6 + 0 + 3 = 9$.
\end{solution}

\begin{solution}{pb:g6:wholes:1}
\textbf{1.} Misalnya $(1 + 4) + (2 + 3) + 5 = 5 + 5 + 5 =
15$.

\textbf{2.} Setiap pasangan bernilai $11$: $1 + 10$, $2 + 9$, $3 + 8$,
$4 + 7$, $5 + 6$. Itu $5$ pasangan (sepuluh bilangan, dua per
pasangan), jadi \omterm{def:g1:addition:def}{jumlahnya} $5 \times 11 = 55$.

\textbf{3.} Pasangannya bernilai $21$ ($1 + 20$, $2 + 19, \dots$), dan
banyaknya $10$: $10 \times 21 = 210$.

\textbf{4.} Pada setiap pasangan, ketika bilangan pertama bertambah
$1$ bilangan kedua berkurang $1$, jadi \omterm{def:g1:addition:def}{jumlahnya} tidak pernah
berubah: $1 + 100 = 2 + 99 = 3 + 98 = \dots = 101$. Seratus bilangan
itu membentuk $50$ pasangan (masing-masing dua bilangan), dengan
pasangan paling dalam $50 + 51 = 101$.

\textbf{5.} $50$ pasangan bernilai $101$:
\[
1 + 2 + \dots + 100 = 50 \times 101 = 5\,050 .
\]

\textbf{6.} Dengan sembilan bilangan, bilangan tengahnya, $5$, tidak
punya pasangan. Di sekelilingnya: $1 + 9 = 2 + 8 = 3 + 7 = 4 + 6 = 10$,
empat pasangan. \omterm{def:g1:addition:def}{Jumlahnya}: $4 \times 10 + 5 = 45$.

\textbf{7.} Setiap kolom bernilai $1 + 9 = 2 + 8 = \dots =
9 + 1 = 10$, dan ada $9$ kolom: kedua baris itu bersama-sama bernilai
$9 \times 10 = 90$. Tetapi kedua baris itu adalah \omterm{def:g1:addition:def}{jumlah} yang
\emph{sama} yang ditulis dua kali, jadi satu salinannya bernilai
setengahnya: $1 + 2 + \dots + 9 = 90 \div 2 = 45$. Tidak ada bilangan
yang tertinggal tanpa pasangan --- pasangannya persis di bawahnya.

\textbf{8.} Untuk $1$ sampai $50$: $50$ kolom bernilai $51$, \omterm{def:g1:addition:def}{jumlah}
gandanya $50 \times 51 = 2\,550$, jadi \omterm{def:g1:addition:def}{jumlahnya} $2\,550 \div 2 = 1\,275$.
Untuk $1$ sampai $200$: $200$ kolom bernilai $201$, \omterm{def:g1:addition:def}{jumlah} gandanya
$40\,200$, \omterm{def:g1:addition:def}{jumlahnya} $20\,100$.

\textbf{9.} Setiap \omterm{def:g3:numbers:evenodd}{bilangan genap} adalah dua kali sebuah bilangan pada
\omterm{def:g1:addition:def}{jumlah} Gauss: $2 = 2 \times 1$, $4 = 2 \times 2$, \dots,
$200 = 2 \times 100$. Jadi seluruh \omterm{def:g1:addition:def}{jumlahnya} dua kali \omterm{def:g1:addition:def}{jumlah} Gauss:
\[
2 + 4 + \dots + 200 = 2 \times 5\,050 = 10\,100 .
\]

\textbf{10.} Bilangan dari $101$ sampai $200$ adalah bilangan dari
$1$ sampai $200$ dengan bilangan $1$ sampai $100$ dibuang:
\[
101 + 102 + \dots + 200 = 20\,100 - 5\,050 = 15\,050 .
\]

\textbf{11.} Bariskan teman-teman itu. Yang pertama berjabat $7$
tangan (satu dengan tiap orang lainnya). Yang kedua sudah menyapa
yang pertama, jadi ia berjabat $6$ tangan yang \emph{baru}; yang
ketiga, $5$; dan seterusnya, yang ketujuh berjabat $1$ dan yang
kedelapan tidak ada lagi yang baru. Setiap jabat tangan terhitung
tepat sekali:
\[
7 + 6 + 5 + 4 + 3 + 2 + 1 = 28 \text{ jabat tangan}
\]
(pasangkan di sekeliling tengahnya: $(7+1) + (6+2) + (5+3) + 4 =
8 + 8 + 8 + 4 = 28$).

\textbf{12.} $1 + 2 + \dots + 12$: dua belas kolom bernilai $13$,
\omterm{def:g1:addition:def}{jumlah} gandanya $12 \times 13 = 156$, jadi
$156 \div 2 = 78$ kaleng.

\textbf{13.} Tangga $1 + 2 + 3 + 4$ persegi, ditambah tangga yang
sama yang dibalik ke bawah, mengisi persis persegi panjang $4$ baris
dan $5$ kolom: setiap baris persegi panjang itu adalah satu kolom
trik dua baris (satu anak tangga yang naik dilengkapi satu anak
tangga yang turun, $1 + 4$, $2 + 3$, $3 + 2$, $4 + 1$ --- selalu
$5$). Jadi dua kali \omterm{def:g1:addition:def}{jumlahnya} adalah $4 \times 5 = 20$ persegi, dan
$1 + 2 + 3 + 4 = 10$: gambarnya \emph{adalah} pertanyaan 7.

\textbf{14.} $1 + 2 + \dots + 12 = 78$ dentang dalam dua belas jam
(pertanyaan 12), jadi $78 \times 2 = 156$ dentang dalam satu hari penuh.

\textbf{15.} Bagilah dengan $60$: $5\,050 = 60 \times 84 + 10$, jadi
$5\,050$ detik adalah $84$ menit dan $10$ detik. Bagilah menitnya
dengan $60$: $84 = 60 \times 1 + 24$, jadi $84$ menit adalah
$1$ jam $24$ menit. Seluruhnya: $1$ jam $24$ menit $10$ detik ---
sedikit \emph{kurang} daripada satu setengah jam ($1$ jam $30$ menit).

\end{solution}
